% SEGMENT OLP-0571-S01
\documentclass[../../../include/open-logic-section]{subfiles}

\begin{document}

\olfileid{sth}{replacement}{ref}
\olsection{மாற்றீடும் பிரதிபலிப்பும்}

\stagesinex{} வழியாக மாற்றீட்டை நியாயப்படுத்தும் நமது கடைசி முயற்சி ஓர்
ஆழமான, அழகிய முடிவுடன் தொடங்குகிறது:\footnote{நினைவூட்டல்: வெளிப்படையாக
வேறுவிதமாகக் கூறப்படாவிட்டால், எல்லா வாய்பாடுகளும் அளவுருக்களைக்
கொண்டிருக்கலாம்.}
%In this, I will use overlining, such as $x_1, \ldots, x_n$, to
%abbreviate ``$x_1, \ldots, x_n$'':
\begin{thm}[பிரதிபலிப்புத் திட்டம்]\ollabel{reflectionschema}
எந்த வாய்பாடு $\phi$ க்கும்:
\[
\forall \alpha \exists \beta > \alpha (\forall x_1 \ldots, x_n \in
V_\beta)(\phi(x_1, \ldots, x_n) \liff \phi^{V_\beta}(x_1, \ldots, x_n))
\]
\end{thm}
\noindent
\olref[sth][replacement][strength]{formularelativization} இல் உள்ளதுபோல்,
$\phi^{V_\beta}$ என்பது $\phi$ இன் ஒவ்வோர் அளவையீட்டுக்குறியையும்
$V_\beta$ என்ற கணத்துக்குக் கட்டுப்படுத்துவதால் கிடைப்பது. ஆகவே
உள்ளுணர்வின்படி பிரதிபலிப்பு பின்வருவதைச் சொல்கிறது: முழுப்
படிநிலைமுறையிலும் $\phi$ மெய்யெனில், படிநிலைமுறையின் தன்னிச்சையாகப் பல
\emph{தொடக்கத் துண்டுகளிலும்} $\phi$ மெய்.

மாற்றீடும் பிரதிபலிப்பும் (பொருத்தமாக முறைப்படுத்தப்பட்ட வடிவங்களில்)
$\Z$ ஐ அடிப்படையாகக் கொண்டால் நிகரானவை என்று \citet{Montague1961} மற்றும்
\citet{Levy1960} காட்டினர்; எனவே இரண்டில் எதைச் சேர்த்தாலும் $\ZF$
கிடைக்கும். (இம்முடிவுகளை \olref[sth][replacement][refproofs]{sec} இல்
நிறுவுகிறோம்.) இந்நிகர்மையைக் கொண்டு, \stagesinex{} வழியாக பிரதிபலிப்பையும்
மாற்றீட்டையும் பின்வருமாறு நியாயப்படுத்தலாம் என்று ஒருவர் நம்பலாம்:
\stagesinex{} கொடுக்கப்பட்டால், படிநிலைமுறை மிக மிக உயரமானதாக இருக்க
வேண்டும்; உண்மையில், அதைப் பற்றி நாம் கூறக்கூடிய எதனாலும் அதன் உயரத்தை
வரம்பிட முடியாத அளவு உயரமானதாக இருக்க வேண்டும். முழுப் படிநிலைமுறையிலும்
எந்த வாக்கியம்~$\phi$ மெய்யாக இருந்தாலும், படிநிலைமுறையின் தன்னிச்சையாகப்
பல தொடக்கத் துண்டுகளிலும் அது மெய்யாக இருக்கும் என்ற சிந்தனையாக இதைப்
புரிந்துகொள்ளலாம். அதுவே பிரதிபலிப்பு.

மீண்டும், மாற்றீட்டுக்கு உள்ளார்ந்த நியாயப்படுத்தலை வழங்கும் உண்மையாகவே
நம்பிக்கையூட்டும் முயற்சிபோல இது தெரிகிறது. ஆனால் இதைப்பற்றி இங்கே கூற
வேண்டியவை மிக அதிகம். அது வெற்றிபெறுகிறதா என்பதை இப்போது நீங்களே
முடிவுசெய்ய வேண்டும்.\footnote{எனினும் \citet[95--100]{Incurvati2020} ஐ
வாசிப்பதன் மூலம் நீங்கள் தொடர விரும்பலாம்.}

\end{document}
